% =============================================================================
% Draft: propagator corrections section for CMS AN-2024/110
% Placement agreed with Giacomo (4 Nov): new Sec. 7.3, directly after 7.2 (LO
% validation with SMEFTsim); current 7.3-7.6 shift by one. Fit results go in
% a new Sec. 9.4 (results_propcorr_local.tex).
%
% Every statement is traced to a source in a "% src:" comment:
%   [G]   SMEFTsim 3.0 practical guide (Brivio), section/eq. numbers
%         (~/Desktop/DrellYan/sources/SMEFTsim_practical_guide.pdf)
%   [UFO] patched model, propagator_hack/SMEFTsim_topU3l_MwScheme_propagatorhack_UFO/
%   [AN]  AN v5 itself
%   [N]   our notes (notes/propagator_correction.md, notes/combine_2Dscans.md)
%
% TODO before pasting into Overleaf: fix the labels marked "TODO label",
% the citation keys, the figure paths and the observations marked TODO.
% =============================================================================

\subsection{EFT parametrization with propagator corrections at LO LHE level}
\label{sec:eft_propcorr}

% src: [G] Sec. 6 intro; Sec. 6.1 option (b) and footnote 16; Eq. (6.9) and the
%      sentence after it.
Dimension-six operators modify the masses and widths of the SM particles, and
for a prediction consistent at $\mathcal{O}(\Lambda^{-2})$ these shifts have to
be included in the propagators~\cite{Brivio:2020onw}.
In the default configuration of \textsc{SMEFTsim}, width corrections are not
included. Since the correction to a propagator is proportional to
$D(q^2)^{-1}$, with $D(q^2) = q^2 - m^2 + i m\Gamma$, it is relevant in the
on-shell region and suppressed far from it~\cite{Brivio:2020onw}: for neutral
DY production this is the region of the Z peak.

% src: [G] Sec. 6.1 ("Adopting a {mW, mZ, GF} input scheme is convenient in this
%      respect, because mass correction effects are entirely avoided"); Eq. (6.10);
%      SMEFTsim implements corrections for Z, W, h, t (Sec. 6.1, first line).
In the $\{m_W, m_Z, G_F\}$ input scheme used in this analysis, mass corrections
are entirely avoided~\cite{Brivio:2020onw}, so only the width correction
$\delta\Gamma_Z$ remains. Of the $Z/\gamma^*$ exchange, only the Z propagator is
corrected. On shell, the relative correction to the propagator is
$-\delta\Gamma_Z/\Gamma_Z^{\mathrm{SM}}$~\cite[Eq.~(6.10)]{Brivio:2020onw}.

\subsubsection{Implementation and generation}

% src: [G] Sec. 6.1 option (a); Eqs. (6.5)-(6.6), (6.9) with delta m = 0;
%      Sec. 6.1 paragraphs on the dummy fields; Appendix A for delta Gamma.
\textsc{SMEFTsim} provides the propagator expanded to linear order in the
shifts~\cite{Brivio:2020onw}. For the Z boson with $\delta m_Z = 0$ this gives
\begin{equation}
  P_Z^{\mu\nu} = P_Z^{\mu\nu,\mathrm{SM}}
  - P_Z^{\mu\nu,\mathrm{SM}}\,\frac{i\, m_Z}{D_Z(q^2)}\,\delta\Gamma_Z ,
  \qquad D_Z(q^2) = q^2 - m_Z^2 + i\, m_Z \Gamma_Z^{\mathrm{SM}} .
  \label{eq:propcorr_linear}
\end{equation}
The correction term is generated as the propagator of a dummy field $Z'$, which
has the same SM interactions as the Z boson and does not enter dimension-six
operators. Vertices with dummy legs are proportional to a flag set by the model
parameter \texttt{linearPropCorrections}. The width shift $\delta\Gamma_Z$ is an internal parameter
of the model, an analytic function of the SM couplings and of the Wilson
coefficients, computed from the Z partial widths~\cite[App.~A]{Brivio:2020onw}.
The correction is switched on in the restriction card with:

% src: production restriction card on LLR (output pasted in the session that
%      fixed the bug); [N] propagator_correction.md Sec. 11-12.
\begin{verbatim}
Block SWITCHES
    1 1.000000e+00 # linearPropCorrections
\end{verbatim}

% src: [G] Sec. 8.4.1 (reweighting with propagator contributions,
%      "launch --rwgt_name=SMEFTsim-cll1-propagator").
Since $\delta\Gamma_Z$ is recomputed from the Wilson coefficients at each
reweighting point, the linear and quadratic templates of
Eq.~(55)
include the propagator correction without changes to the reweighting procedure.

% src: [G] Sec. 8.3.1 (NP==1 NPprop=0 vs NP=0 NPprop... examples; "propagator
%      corrections from different operators and EFT orders cannot be
%      disentangled"; NPprop expansion_order >= 2, default 0);
%      [UFO] coupling_orders.py (NPall, NPprop expansion_order 99),
%      couplings.py (dim-6 vertices: NPall 2; dummy vertices: NPprop 1, NPall 1).
In \textsc{SMEFTsim}, vertex and propagator insertions are selected by separate
interaction orders, \texttt{NP} and \texttt{NPprop}, and the latter is disabled
by default~\cite{Brivio:2020onw}. To generate both at the same order in the EFT
expansion in a single sample, we added to the model an interaction order
\texttt{NPall}: each dimension-six vertex
counts as two units and each $Z'$ vertex as one, so a single $Z'$ propagator
counts as one dimension-six insertion. The process is generated with:

% src: proc card DYSMEFTMll50_120_propcorr_proc_card.dat (genproductions
%      cards/DY_SMEFT_Gridpacks/), relevant lines only.
\begin{verbatim}
import model SMEFTsim_topU3l_MwScheme_UFO-all_massless
generate p p > mu+ mu- QCD=0 SMHLOOP=0 NPall<=2 NPprop^2<=2
\end{verbatim}

% src: [G] Sec. 8.2.2 (XX=n at amplitude level, XX^2 at squared-amplitude level);
%      Sec. 8.3.1 (NPprop^2==2 -> interference prop. to delta Gamma_Z;
%      NPprop==2 -> quadratic, prop. to delta Gamma_Z^2).
\noindent At amplitude level, \texttt{NPall<=2} allows either one dimension-six
vertex or one $Z'$ propagator. At squared-amplitude level, \texttt{NPprop\^{}2<=2}
keeps the interference proportional to $\delta\Gamma_Z$ and removes the term
proportional to $(\delta\Gamma_Z)^2$~\cite{Brivio:2020onw}. Without propagator
corrections the same process is generated with \texttt{NP<=1}.

% src: [N] propagator_correction.md Sec. 14, 16, 17; combine_2Dscans.md (406 launch
%      blocks in the reweight card).
Gridpacks with propagator corrections are produced in the same seven
$m_{\ell\ell}$ bins as those without, with the same restriction card and the same
406 reweighting points (SM point, $\pm1$ for each operator, and all operator
pairs). For the validation at LHE level, 100k events per bin are generated.
% NanoAOD production deliberately not described here: the samples used for the
% results are Giacomo's large production (ask him how he wants it documented).

% Optional, to be decided with Giacomo: mention the earlier invalid production.
% A first production was found to have \texttt{linearPropCorrections} set to zero
% in the model used on the worker nodes, so that the $Z'$ diagrams did not
% contribute; it was discarded and all samples were regenerated.

% src: propcorr: XSECUP of the LHE <init> block, first seed file per bin
%      (spot-checked on a second seed file); baseline: samples.json entries
%      DYSMEFTsim_LO_mll_<bin>; ratios computed from these numbers
%      (50-120: 1088.54/1188.67 = 0.916; others 0.996-1.003, see open_questions.md).
%      Interpretation: [G] Sec. 6 (correction relevant on shell).
At the generation point, with all operators at unitary Wilson coefficient, the
propagator correction changes the cross section by $-8.4\%$ in the
$50 < m_{\ell\ell} < 120\GeV$ bin and by at most 0.4\% in the other bins,
consistently with a correction concentrated on shell.

% src: [UFO] parameters.py: dWZ contains cHj1, cHj3, cHQ1, cHQ3, cHu, cHd, cHbq,
%      cHl1, cHl3, cHe, cHWB, cbWRe, cbBRe directly, and dg1, dgw, which depend on
%      dGf = (2 cHl3 - cll1) v^2/L^2 and dMZ2 = (cHDD/2 + 2 cHWB cth sth) v^2/L^2;
%      [G] Table 4 (topU3l row of Delta G_F, Delta m_Z^2), App. A Eqs. (A.2)-(A.5),
%      (A.9)-(A.13); cross-checked with automated diagram inspection
%      (auto_detect_operators_propcorr.py, [N] propagator_correction.md Sec. 9).
Of the operators affecting neutral DY production
(Table~20),
15 enter $\delta\Gamma_Z$. They can be read off the expression of
$\delta\Gamma_Z$ in the model, and were cross-checked by inspecting the generated
diagrams for $Z'$ insertions~\cite{EFT2Obs}:
\begin{itemize}
  \item through the Z couplings to quarks: $\mathcal{O}_{Hq}^{(1)}$,
    $\mathcal{O}_{Hq}^{(3)}$, $\mathcal{O}_{HQ}^{(1)}$, $\mathcal{O}_{HQ}^{(3)}$,
    $\mathcal{O}_{Hu}$, $\mathcal{O}_{Hd}$, $\mathcal{O}_{Hb}$, and the dipoles
    $\mathcal{O}_{bW}$, $\mathcal{O}_{bB}$;
  \item through the Z couplings to leptons: $\mathcal{O}_{H\ell}^{(1)}$,
    $\mathcal{O}_{H\ell}^{(3)}$, $\mathcal{O}_{He}$;
  \item through the input parameters, shifting all Z couplings at once:
    $\mathcal{O}_{HD}$ and $\mathcal{O}_{HWB}$ via $m_Z$, and
    $\mathcal{O}_{H\ell}^{(3)}$ and $\mathcal{O}_{\ell\ell}^{(1)}$ via
    $G_F$~\cite[Table~4]{Brivio:2020onw}.
\end{itemize}
% src: [UFO] parameters.py dWZ (no four-fermion coefficient appears).
The remaining operators are four-fermion contact interactions; they do not
appear in $\delta\Gamma_Z$, so their templates are not affected.

% src: [UFO] dWZ contains cHQ1 and cHQ3 only in the sum; [G] Eq. (A.5)
%      (Delta^L_d = C_Hq^(1) + C_Hq^(3) for down-type quarks); identical
%      event weights observed in the LHE files ([N] propagator_correction.md Sec. 15).
%      dipoles: [UFO] dWZ, cbWRe/cbBRe only multiplied by MB; [G] Eq. (A.12);
%      [AN] Sec. 7.2 (vanishing ObB/ObW).
Two features follow from these expressions. $\mathcal{O}_{HQ}^{(1)}$ and
$\mathcal{O}_{HQ}^{(3)}$ enter only through their sum, both in $\delta\Gamma_Z$ and
in the Z coupling to b quarks~\cite[Eq.~(A.5)]{Brivio:2020onw}, so their templates
are identical event by event and only the sum can be constrained. The dipole
operators enter $\delta\Gamma_Z$ only through terms proportional to the b-quark
mass~\cite[Eq.~(A.12)]{Brivio:2020onw}, in line with their vanishing vertex
contributions in Section~7.2.


\subsubsection{Linear and quadratic templates}

% src: [AN] Eq. 55 (eq:eft_parametrization_1op, N(k) from w(0), w(+-1)); Sec. 7.2 (comparison per operator); compare_ops.py
%      (w_lin = (w_p1 - w_m1)/2, w_quad = (w_p1 + w_m1)/2 - w_SM; defaults 80-102 GeV,
%      0.5 GeV bins; --no-normalize; colours); cache built with build_cache_propcorr.py
%      ([N] propagator_correction.md Sec. 6, 15).
As in Section~7.2,
the linear and quadratic templates are derived with
Eq.~(55),
with and without propagator corrections, and compared in the range
$80 < m_{\ell\ell} < 102\GeV$
(Figs.~\ref{fig:propcorr_lin} and~\ref{fig:propcorr_quad}).

% src: [G] Sec. 6.1 (a) (linear dependence on the Wilson coefficients at amplitude
%      level); Sec. 8.3.1 (Z' diagrams interfering with the SM: "pure SM-L6
%      interference (prop. to delta Gamma_Z)"); [UFO] dWZ is linear in the Wilson
%      coefficients, so this term is odd in C and lands in w_lin.
%      Observations: compare_<op>_lin_raw.pdf (5 Nov, LHE/propcorr/plots_AN, default
%      100k baseline cache). Peak values (pb/bin), baseline -> propcorr:
%      cHl3 -10 -> +4.7; cHj3 +12.5 -> -1.7; cHQ1/3 +0.6 -> -3.0; cHu +2.4 -> -0.5;
%      cHWB +5.2 -> ~+1; cll1 +11 -> ~0; cHj1 +1.9 -> ~+0.7; cHl1 ratio ~1.1;
%      cHe ratio ~0.8; cHbq -0.05 -> +0.6; cbWRe -0.09 -> +0.43; cbBRe -0.05 -> +0.23;
%      cHDD, cHd: statistically limited. Locality: [G] Sec. 6 (correction prop. to
%      1/D(q^2), relevant on shell).
The linear term now also contains the interference of the $Z'$ diagrams with
the SM amplitude, proportional to $\delta\Gamma_Z$~\cite{Brivio:2020onw}. Its
effect is confined to the Z pole, where it is often larger than the vertex
contribution and can change the sign of the linear term, as for
$\mathcal{O}_{H\ell}^{(3)}$ and $\mathcal{O}_{HQ}^{(1,3)}$.
% src: compare_<op>_quad_raw.pdf (5 Nov). Peak values (pb/bin), baseline -> propcorr:
%      cll1 +0.34 -> -0.32; cHl3 +0.55 -> -0.27; cHj3 +0.49 -> -0.43;
%      cHQ1/3 0.022 -> 0.010; cHWB, cHe ~10-20% lower; cHl1 ~15% higher;
%      cHDD, cHj1, cHu, cHd compatible; cHbq, cbWRe, cbBRe, cQe statistically limited.
%      Interpretation: amplitudes allowed by NPall<=2 are SM, one dim-6 vertex (prop.
%      to C) or one Z' propagator (prop. to delta Gamma_Z, linear in C, [UFO] dWZ);
%      NPprop^2<=2 removes |Z'|^2 ([G] Secs. 8.2.2, 8.3.1); so the extra C^2 term is
%      the dim-6-vertex x Z' interference, which can be negative.
The quadratic templates change as well, and can become negative at the pole, as
for $\mathcal{O}_{\ell\ell}^{(1)}$: they now include the interference between
the diagrams with a dimension-six vertex and those with a $Z'$ propagator.
% src: [UFO] parameters.py dWZ (no four-fermion coefficient); compare_cQe_*_raw.pdf.
The templates of four-fermion operators, which do not enter $\delta\Gamma_Z$,
are unchanged (Fig.~\ref{fig:propcorr_control}).

\begin{figure}[htbp]
  \centering
  % from run_compare_all_ops.sh --no-normalize (compare_<op>_lin_raw.pdf)
  \includegraphics[width=0.24\textwidth]{figures/propcorr/compare_cHDD_lin_raw.pdf}%
  \hfill
  \includegraphics[width=0.24\textwidth]{figures/propcorr/compare_cHWB_lin_raw.pdf}%
  \hfill
  \includegraphics[width=0.24\textwidth]{figures/propcorr/compare_cll1_lin_raw.pdf}%
  \hfill
  \includegraphics[width=0.24\textwidth]{figures/propcorr/compare_cHl1_lin_raw.pdf}%
  \\
  \includegraphics[width=0.24\textwidth]{figures/propcorr/compare_cHl3_lin_raw.pdf}%
  \hfill
  \includegraphics[width=0.24\textwidth]{figures/propcorr/compare_cHe_lin_raw.pdf}%
  \hfill
  \includegraphics[width=0.24\textwidth]{figures/propcorr/compare_cHj1_lin_raw.pdf}%
  \hfill
  \includegraphics[width=0.24\textwidth]{figures/propcorr/compare_cHj3_lin_raw.pdf}%
  \\
  \includegraphics[width=0.24\textwidth]{figures/propcorr/compare_cHQ1_lin_raw.pdf}%
  \hfill
  \includegraphics[width=0.24\textwidth]{figures/propcorr/compare_cHQ3_lin_raw.pdf}%
  \hfill
  \includegraphics[width=0.24\textwidth]{figures/propcorr/compare_cHu_lin_raw.pdf}%
  \hfill
  \includegraphics[width=0.24\textwidth]{figures/propcorr/compare_cHd_lin_raw.pdf}%
  \\
  \includegraphics[width=0.24\textwidth]{figures/propcorr/compare_cHbq_lin_raw.pdf}%
  \hfill
  \includegraphics[width=0.24\textwidth]{figures/propcorr/compare_cbWRe_lin_raw.pdf}%
  \hfill
  \includegraphics[width=0.24\textwidth]{figures/propcorr/compare_cbBRe_lin_raw.pdf}%
  \hfill
  \caption{Linear ($\Lambda^{-2}$) templates in the Z-peak region at LHE level for
    the 15 operators entering $\delta\Gamma_Z$, without (blue) and with (red)
    propagator corrections. The bottom panels show the ratio of the two.}
  \label{fig:propcorr_lin}
\end{figure}

\begin{figure}[htbp]
  \centering
  \includegraphics[width=0.24\textwidth]{figures/propcorr/compare_cHDD_quad_raw.pdf}%
  \hfill
  \includegraphics[width=0.24\textwidth]{figures/propcorr/compare_cHWB_quad_raw.pdf}%
  \hfill
  \includegraphics[width=0.24\textwidth]{figures/propcorr/compare_cll1_quad_raw.pdf}%
  \hfill
  \includegraphics[width=0.24\textwidth]{figures/propcorr/compare_cHl1_quad_raw.pdf}%
  \\
  \includegraphics[width=0.24\textwidth]{figures/propcorr/compare_cHl3_quad_raw.pdf}%
  \hfill
  \includegraphics[width=0.24\textwidth]{figures/propcorr/compare_cHe_quad_raw.pdf}%
  \hfill
  \includegraphics[width=0.24\textwidth]{figures/propcorr/compare_cHj1_quad_raw.pdf}%
  \hfill
  \includegraphics[width=0.24\textwidth]{figures/propcorr/compare_cHj3_quad_raw.pdf}%
  \\
  \includegraphics[width=0.24\textwidth]{figures/propcorr/compare_cHQ1_quad_raw.pdf}%
  \hfill
  \includegraphics[width=0.24\textwidth]{figures/propcorr/compare_cHQ3_quad_raw.pdf}%
  \hfill
  \includegraphics[width=0.24\textwidth]{figures/propcorr/compare_cHu_quad_raw.pdf}%
  \hfill
  \includegraphics[width=0.24\textwidth]{figures/propcorr/compare_cHd_quad_raw.pdf}%
  \\
  \includegraphics[width=0.24\textwidth]{figures/propcorr/compare_cHbq_quad_raw.pdf}%
  \hfill
  \includegraphics[width=0.24\textwidth]{figures/propcorr/compare_cbWRe_quad_raw.pdf}%
  \hfill
  \includegraphics[width=0.24\textwidth]{figures/propcorr/compare_cbBRe_quad_raw.pdf}%
  \hfill
  \caption{Quadratic ($\Lambda^{-4}$) templates in the Z-peak region at LHE level
    for the 15 operators entering $\delta\Gamma_Z$, without (blue) and with (red)
    propagator corrections. The bottom panels show the ratio of the two.}
  \label{fig:propcorr_quad}
\end{figure}

\begin{figure}[htbp]
  \centering
  \includegraphics[width=0.32\textwidth]{figures/propcorr/compare_cQe_lin_raw.pdf}
  \includegraphics[width=0.32\textwidth]{figures/propcorr/compare_cQe_quad_raw.pdf}
  \caption{Linear (left) and quadratic (right) templates in the Z-peak region at
    LHE level for a four-fermion operator not entering $\delta\Gamma_Z$, without
    (blue) and with (red) propagator corrections. The bottom panels show the
    ratio of the two.}
  \label{fig:propcorr_control}
\end{figure}
